\documentclass[11pt,a4paper]{report}
\usepackage{mypka-document}

\title{A Reproducible Research Document}
\author{Author Name}
\date{\today}

\begin{document}
\maketitle
\tableofcontents

\chapter{Question and contribution}
The document states the question, evidence, result, and limitation before detail. \citet{example2025} supplies a narrative citation. Related work supplies a parenthetical multi-work example \citep{example2025,companion2024}.

\section{Lists}
Lists use generous but controlled spacing.
\begin{itemize}
  \item Lead with the substantive claim.
  \item Separate evidence from inference.
  \item State the principal limitation.
\end{itemize}
\begin{enumerate}
  \item Define the question.
  \item Present the design.
  \item Interpret the result.
\end{enumerate}

\section{Native table}
\begin{table}[H]
  \centering
  \caption{Illustrative estimates.}
  \label{tab:estimates}
    \begin{tabular}{lcc}
      \toprule
      Specification & {Estimate} & {Std. error} \\
      \midrule
      Baseline & 0.42 & 0.08 \\
      Robustness & 0.39 & 0.09 \\
      \bottomrule
    \end{tabular}

    \parbox{0.72\textwidth}{\footnotesize\textit{Notes:} Values are illustrative. The table contains no empirical claim.}
\end{table}

\section{TikZ plot and diagram}
\begin{figure}[H]
  \centering
  \begin{tikzpicture}[x=1.25cm,y=1.25cm]
\draw[->] (0,0) -- (5.7,0) node[right] {Period};
\draw[->] (0,0) -- (0,3.2) node[above] {Index (units)};
\foreach \x in {1,...,5} \draw (\x,0.06) -- (\x,-0.06) node[below] {\x};
\foreach \y in {1,2,3} \draw (0.06,\y) -- (-0.06,\y) node[left] {\y};
\draw[mypkalink,very thick,mark=*] plot coordinates {(1,1.0) (2,1.4) (3,1.8) (4,2.3) (5,2.7)};
\node[anchor=west,mypkalink] at (0.35,2.85) {Illustrative series};
\end{tikzpicture}

  \caption{Illustrative series rendered with PGFPlots.}
  \label{fig:plot}
\end{figure}
\begin{figure}[H]
  \centering
  \begin{tikzpicture}[
  node distance=1.8cm,
  box/.style={draw=mypkalink,rounded corners,align=center,minimum width=3.1cm,minimum height=0.9cm},
  flow/.style={-{Latex[length=3mm]},thick,mypkalink}
]
\node[box] (question) {Research question};
\node[box,right=of question] (design) {Design or model};
\node[box,right=of design] (result) {Result and limitation};
\draw[flow] (question) -- (design);
\draw[flow] (design) -- (result);
\end{tikzpicture}

  \caption{Illustrative research flow rendered with TikZ.}
  \label{fig:diagram}
\end{figure}

\chapter{Conclusion}
The source demonstrates cross-references to Table~\ref{tab:estimates}, Figure~\ref{fig:plot}, and Figure~\ref{fig:diagram}, together with a reproducible bibliography.

\bibliographystyle{apalike}
\bibliography{references}
\end{document}
